\section{BMW's Quality Process and Causer Attribution}\label{sec:fund_bmw}

BMW's quality management process encompasses the detection and
documentation of product defects, the containment of their effects, the
attribution of responsibility, and the initiation of appropriate
problem-solving activities. In the context of this thesis, the central
process-level decision is whether a detected defect is attributable to a
supplier or to an internal unit, such as assembly, logistics, or another
production-related function. This early attribution determines the
subsequent investigation and complaint route. The remainder of this
section describes the systems, data, and roles involved in this process,
from initial defect detection (Section~\ref{sec:fund_bmw_detection}) to
the supplier complaint and 8D process (Section~\ref{sec:fund_bmw_8d}).
Figure~\ref{fig:bmw_process_overview} summarizes this process at a high
level before each stage is described in detail.

\begin{minipage}{\textwidth}
\centering
\begin{tikzpicture}[
  scale=0.82, transform shape,
  node distance=0.45cm and 1.2cm,
  process/.style={
    rectangle, rounded corners, draw=black, thick,
    minimum width=7cm, minimum height=1.1cm,
    align=center, font=\small, inner sep=4pt
  },
  decision/.style={
    diamond, draw=black, thick, aspect=3,
    minimum width=5.4cm, minimum height=1.8cm,
    align=center, font=\small, inner sep=2pt
  },
  outcome/.style={
    rectangle, draw=black, thick, double, double distance=1pt,
    minimum width=4.6cm, minimum height=1.2cm,
    align=center, font=\small, inner sep=4pt
  },
  lbl/.style={font=\scriptsize},
  arrow/.style={-{Latex[length=2mm]}, thick}
]

\node[process] (detect) {Defect Detected\\\footnotesize(IPS-Q via MQC)};
\node[process, below=of detect] (describe)
  {Defect Described via VFC\\\footnotesize(Defect Type, Place, Location)};
\node[process, below=of describe] (qi)
  {Quality Incident (QI) Created\\\footnotesize(if formal investigation required)};
\node[process, below=of qi] (evidence)
  {Historical Evidence Reviewed\\\footnotesize(exact-match lookup, inspection, sorting, rework)};
\node[decision, below=0.8cm of evidence] (decision) {Causer Decision};
\node[outcome, below left=1.6cm and -0.3cm of decision] (supplier)
  {Supplier Complaint\\\footnotesize\& 8D Process};
\node[outcome, below right=1.6cm and -0.3cm of decision] (internal)
  {Internal Investigation\\\footnotesize(Assembly, Logistics, Parts Quality)};

\draw[arrow] (detect) -- (describe);
\draw[arrow] (describe) -- (qi);
\draw[arrow] (qi) -- (evidence);
\draw[arrow] (evidence) -- (decision);
\draw[arrow] (decision.south) -- ++(0,-0.45) -| node[lbl, pos=0.2, above] {supplier} (supplier.north);
\draw[arrow] (decision.south) -- ++(0,-0.45) -| node[lbl, pos=0.2, above] {internal} (internal.north);

\end{tikzpicture}
\captionsetup{type=figure}
\captionof{figure}{Overview of BMW's current quality-incident and
  causer-attribution process, synthesized from
  Sections~\ref{sec:fund_bmw_detection}--\ref{sec:fund_bmw_8d}.}
\label{fig:bmw_process_overview}
\end{minipage}
\vspace{\baselineskip}

\subsection{Defect Detection and Recording}\label{sec:fund_bmw_detection}

The process begins when a defect is detected during production,
inspection, or rework. A defect may concern, for example, the appearance,
function, condition, installation, or completeness of a component. The
\textbf{International Production System for Quality (IPS-Q)} provides the
production-side context for defects and rework records. Internal material
identifies \textbf{Modern Quality Capturing (MQC)} as the mobile IPS-Q
application used for defect recording and initial causer allocation.

A defect record may contain the affected material, supplier, plant,
vehicle, detection time, a textual problem description, and standardized
defect characteristics. The IPS-Q-related data fields identified for this
process include:

\begin{itemize}
  \item the material number;
  \item the supplier number;
  \item the plant;
  \item the timestamp;
  \item the defect description; and
  \item vehicle identification information.
\end{itemize}

The quality of the initial textual description is important, because it
determines whether comparable historical cases can be found and correctly
interpreted. The QMT front-desk process explicitly asks whether the
problem description is complete and understandable \emph{before}
considering repeat occurrence, causer attribution, rework, potential root
causes, and short-term actions. A description such as ``door damaged''
provides little information for historical comparison, whereas a
description such as ``damage to the ambient lighting area of the front
trim'' identifies the affected component and physical position more
precisely.

\subsection{Standardized Fault Coding}\label{sec:fund_bmw_vfc}

To describe defects consistently, BMW uses a standardized fault coding
scheme referred to as the \textbf{Vehicle Fault Code (VFC; German:
\textit{Vereinheitlichte Fehlercodierung})}.
The VFC provides a standardized structure for documenting product-related
defects across organizational units, systems, and processes. For this
thesis, three VFC elements are particularly relevant: the \textbf{defect
type}, which describes the observed form of deviation and effect; the
\textbf{defect place}, which identifies the affected component or
functional area; and the \textbf{defect location}, which provides a more
precise physical or positional localization. Internal data descriptions
contain corresponding fields for defect type, defect place, and defect
location, which are used together with material, supplier, plant, and
vehicle information to identify defect patterns.

Table~\ref{tab:vfc_elements} illustrates these three elements, together
with the shared material key, using the door-trim example introduced
above.

\begin{minipage}{\textwidth}
\centering
\footnotesize
\renewcommand{\arraystretch}{0.9}
\captionsetup{type=table,skip=3pt}
\captionof{table}{VFC elements and the shared material key, illustrated using the door-trim damage example.}
\label{tab:vfc_elements}
\begin{tabular}{lp{5cm}p{3.6cm}}
  \toprule
  \textbf{Element} & \textbf{Description} & \textbf{Example Value} \\
  \midrule
  Defect Type &
    Observed form of deviation and effect &
    Damage \\
  Defect Place &
    Affected component or functional area &
    Door trim, ambient lighting \\
  Defect Location &
    Precise physical or positional localization &
    Front right \\
  Material (shared key) &
    Identifies the affected part &
    Corresponding door-trim part number \\
  \bottomrule
\end{tabular}
\end{minipage}
\vspace{\baselineskip}

In this thesis, the ordered combination of defect type, defect place, and
defect location is referred to as the \textbf{defect signature}:
\[
  D = (T, P, L),
\]
where $T$ denotes the defect type, $P$ the defect place, and $L$ the
defect location. The defect signature provides a common basis for
retrieving related records: the knowledge graph construction described in
Chapter~\ref{ch:problem} likewise treats defect type, defect place,
defect location, and material as business keys for linking historical
quality incidents and associated rework evidence.

A shared defect signature, however, does not by itself establish an
identical causer. The damage in the door-trim example could have been
present when the component arrived from the supplier, or it could have
occurred during internal transport, assembly handling, or rework. The VFC
standard describes what was observed; it does not independently determine
how the defect arose or who caused it.

\subsection{Quality Incident Management}\label{sec:fund_bmw_qi}

When a defect requires formal investigation, it is documented as a
\textbf{Quality Incident (QI)}. A QI represents the investigation of a
quality problem rather than only an individual defect occurrence. Internal
process documentation defines a QI as the formal documentation of a defect
detected during production, followed by an investigation to identify the
responsible unit. The relevant process documentation associates QI
creation, messages, and tasks with QCat and Qcube, systems that also
support activities for planning and documenting containment, sorting,
rework, and inspection.

The distinction between a \textbf{defect occurrence} and a QI is
important:

\begin{itemize}
  \item A \emph{defect occurrence} describes a specific observation on a
    vehicle or component.
  \item A \emph{QI} provides the broader investigation context in which
    one or more observations are examined.
\end{itemize}

For example, the same door-trim damage may be detected on several
vehicles; each vehicle may generate a separate inspection or rework
record. These records may be investigated under one QI if they are
considered manifestations of the same quality problem. Conversely, two
defects with the same VFC description may belong to different QIs if they
involve different materials, suppliers, plants, or process conditions. The
QI therefore serves as the formal context in which a defect, its possible
recurrence, and its potentially responsible units are assessed.

\subsection{Process-Level Causer Attribution}\label{sec:fund_bmw_causer}

\textbf{Process-level causer attribution} refers to the identification of
the organizational or process-level unit most likely responsible for a
quality defect. The initial decision generally distinguishes between a
supplier-related issue and an internally caused issue involving assembly,
logistics, or another relevant unit. The QMT front-desk process explicitly
includes the question of who the causer for the issue is, with assembly,
parts quality, and logistics listed as possible areas to be considered.

Causer attribution must be distinguished from \textbf{root-cause
identification}. Causer attribution determines the organizational or
process-level area to which an incident should initially be directed;
root-cause identification determines the precise technical or procedural
mechanism that produced the defect. For example, assembly may initially be
identified as the likely causer before the technical root cause is found
to be an unsuitable installation tool or an incorrect handling sequence.
Similarly, a supplier may be identified as the likely responsible party
before a manufacturing or packaging deviation is confirmed as the
technical root cause.

The QMT front-desk assessment includes more than the defect
classification. The documented business questions address:

\begin{itemize}
  \item whether the problem description is complete;
  \item whether the issue has occurred previously (repeat occurrence);
  \item who the likely causer is;
  \item whether the rework has been fully and correctly recorded;
  \item which potential root causes and short-term actions are known;
  \item whether the issue is isolated or part of a series; and
  \item whether other plants are, or may be, affected.
\end{itemize}

The result is an early process decision based on the information
available at the time. It guides the subsequent investigation but may be
revised as additional evidence becomes available.

\subsection{Historical Lookup and Repeat-Problem Identification}\label{sec:fund_bmw_lookup}

The current investigation process combines structured historical lookup
with expert assessment. Internal QMT training material describes a
repeat-failure search that queries three different combinations of
supplier number, part number, and VFC attributes.
% TODO: confirm the exact name of the lookup tool/function described in
% the QMT training material -- unclear in the source recording.
An empty result list indicates that no comparable QI was found using the
active search configuration; where a QI with the same defect pattern
already exists, the new QI can be linked to the original case and marked
as a repeat failure.

At a high level, the current process can be summarized as follows:

\begin{enumerate}
  \item A defect is recorded in IPS-Q through MQC.
  \item The defect is described using the VFC and relevant product and
    process information.
  \item A QI is created when formal investigation is required.
  \item Historical incidents are searched using configured combinations of
    supplier, part, and VFC attributes.
  \item Related inspection, sorting, and rework information is reviewed.
  \item The incident is assigned to a supplier or an internal unit.
  \item The responsible expert validates the assignment and initiates the
    corresponding follow-up process.
\end{enumerate}

This sequence is a synthesis of the documented system landscape and
front-desk business questions; the exact execution may vary according to
the type and context of the incident. For example, an investigator
assessing door-trim damage may retrieve incidents involving the same
material and VFC description. If comparable incidents were previously
assigned to the supplier, they provide historical support for supplier
attribution. Nevertheless, their applicability to the current case must be
evaluated using the supplier, plant, detection point, production stage,
and other available context.

The current lookup is transparent, because the filter attributes are
explicit. Its coverage is nevertheless limited by the defined search
combinations and the completeness of the recorded information: a relevant
historical case may not be returned if it differs in material, supplier,
fault type, or another required attribute. The thesis proposal identifies
this exact-match limitation as a central weakness of the existing
approach.

\subsection{Inspection, Sorting, and Rework}\label{sec:fund_bmw_isr}

Inspection, sorting, and rework are related quality activities, but they
serve different purposes and produce different forms of information.
\textbf{Inspection} evaluates a product's quality characteristics and
records the context in which a defect was identified; the inspection
information considered in this thesis includes the station, the
production stage, the relevant quality check, and the test result.
\textbf{Sorting} examines a defined population of parts or vehicles, and
it distinguishes sorting performed by a supplier, internally by BMW, or
through a mixed form involving BMW personnel, a quality service provider,
and the supplier. \textbf{Rework} documents how a detected defect was
corrected; the relevant records may contain the defect type, the repair
action (referred to as the rework activity), the affected component, the
plant, and the responsible department.

Across the different systems used by QMT, these activities are
represented through \textbf{Sorting, Rework, and Inspection (SRI) orders}.
The documented SRI process includes order creation, instructions, and the
recording of inspection results as OK or NOK (not OK).

These activities provide complementary perspectives: inspection records
where and under which conditions a defect is detected, sorting records the
result of examining a defined population, and rework records how the
defect was corrected. These perspectives must not be confused with
causation. The place at which a defect is detected is not necessarily the
place at which it was created; likewise, the unit that performs the repair
is not automatically the causer. For example, door-trim damage detected at
an assembly inspection point may already have existed before the
component entered assembly; conversely, a conforming component may be
damaged during installation and subsequently reworked at another station.
Detection, correction, and causation therefore remain separate aspects of
the investigation.

\subsection{Roles in the Quality Process}\label{sec:fund_bmw_roles}

Several roles contribute to defect detection, investigation, containment,
and problem-solving:

\begin{description}[%
  style=nextline,
  leftmargin=0pt,
  labelsep=0pt,
  font=\normalfont\bfseries,
  itemsep=\baselineskip,
  topsep=\smallskipamount
]
  \item[Production and Assembly]
    Identify defects in the manufacturing environment and provide the
    initial defect and process context. MQC is positioned in the internal
    ecosystem as the application for recording a defect and performing an
    initial allocation to a likely causer.

  \item[Quality Engineers and QMT Front-Desk Personnel]
    Examine the problem description, compare the issue with previous
    cases, assess a likely causer, and determine the appropriate follow-up
    process. Their documented business questions include repeat
    occurrence, causer identification, rework completeness, potential
    root causes, short-term actions, and affected plants.

  \item[Internal Units (Assembly, Logistics, Parts Quality, and Related Units)]
    May receive incidents for further investigation when internal
    responsibility is suspected. These units appear as potential causer
    categories in the processing paths documented in the QMT front-desk
    material.

  \item[Suppliers]
    Participate in the complaint and problem-solving process when
    supplier responsibility is suspected. A dedicated supplier portal,
    qWin, supports the exchange and management of supplier-related 8D
    information, including requested 8D reports. External quality
    service providers may perform sorting activities,
    while mixed sorting additionally involves the supplier.
\end{description}

\subsection{Supplier Complaints and the 8D Process}\label{sec:fund_bmw_8d}

If a defect is considered supplier-related, the incident can proceed to a
supplier complaint process. The documented system landscape connects the
quality notification to \textbf{8D (eight-disciplines) reporting}.

The 8D method is a structured problem-solving approach used to identify,
correct, and prevent recurring product or process problems. It progresses
from team formation and problem definition, through containment and
root-cause analysis, to permanent corrective actions and the prevention of
recurrence.

The 8D method originated in the automotive industry as a standardized
framework for supplier and internal problem-solving and is applied within
BMW's quality management process whenever a defect is confirmed as
supplier-related. Although the disciplines are numbered sequentially, they
are not always executed strictly one after another in practice: D3
(interim containment) is typically initiated in parallel with, or even
before, the completion of D4 (root-cause analysis), since protecting the
customer and the production process cannot wait until the root cause is
fully understood.

In practice, each QI is assigned a priority level, and the eight
disciplines are governed by corresponding target timelines that differ
between high- and low-priority cases.
% TODO: confirm the exact priority levels/terminology and timeline
% structure with the industry supervisor -- unclear in the source
% recording.

The cross-functional team formed under D1 typically draws on the roles
introduced in Section~\ref{sec:fund_bmw_roles}: the responsible quality
engineer coordinates the investigation, while representatives from
production, assembly, and, where relevant, the supplier contribute the
domain-specific knowledge needed to describe the problem, contain it, and
identify its root cause.

The 8D reports generated through this process also complement the
historical lookup described in Section~\ref{sec:fund_bmw_lookup}: whereas
that lookup retrieves comparable past QIs based on supplier, part, and
VFC attributes, an attached 8D report reveals not only that a comparable
incident existed, but how it was investigated and resolved, and through
which discipline-specific findings.

Each discipline is documented with a responsible owner, supporting
evidence, and a completion date. This staged documentation means that an
8D report's evidentiary value depends on how far the investigation had
progressed at the time it is reviewed: a report that is still open may
contain only a preliminary problem description and containment action,
whereas a closed report additionally documents a verified root cause and
validated corrective actions.

For the purposes of this thesis, this discipline structure matters because
an 8D report is not a single, undifferentiated text block: each discipline
corresponds to a distinct, dated piece of evidence, such as a problem
description, a containment action, or a verified root cause, that can in
principle be linked to the underlying QI as a separate item of evidence
rather than treated as one opaque document. Table~\ref{tab:8d_disciplines}
summarizes the eight disciplines reviewed for this thesis.

\begin{minipage}{\textwidth}
\centering
\footnotesize
\captionsetup{type=table}
\captionof{table}{The eight disciplines (8D) of the problem-solving process.}
\label{tab:8d_disciplines}
\renewcommand{\arraystretch}{1.05}
\begin{tabular}{lp{3.9cm}p{6.8cm}}
  \toprule
  \textbf{Discipline} & \textbf{Title} & \textbf{Purpose} \\
  \midrule
  D1 &
    Establish the Team &
    Form a cross-functional team with the product and process knowledge
    required to investigate and resolve the problem. \\
  \addlinespace
  D2 &
    Define and Describe the Problem &
    Describe the problem clearly and, where possible, quantitatively. A
    description should establish what occurred, where and when it
    occurred, who or what was affected, and the extent of the problem. \\
  \addlinespace
  D3 &
    Develop and Implement Interim Containment Actions &
    Introduce temporary measures to isolate the problem and protect the
    customer or production process while the root cause is being
    investigated. \\
  \addlinespace
  D4 &
    Determine and Verify Root Causes &
    Identify and verify the causes that explain why the problem occurred.
    This discipline also considers why the defect was not detected
    earlier. Methods such as the 5~Whys and cause-and-effect diagrams may
    support the analysis. \\
  \addlinespace
  D5 &
    Select and Verify Permanent Corrective Actions &
    Identify permanent actions that address the verified root cause, and
    confirm that the selected actions are capable of resolving the
    problem. \\
  \addlinespace
  D6 &
    Implement and Validate Corrective Actions &
    Implement the selected permanent corrective actions and verify their
    effectiveness using appropriate evidence. \\
  \addlinespace
  D7 &
    Prevent Recurrence &
    Modify relevant systems, procedures, practices, or controls to
    prevent the same or similar problems from recurring. \\
  \addlinespace
  D8 &
    Recognize the Team and Close the Process &
    Recognize the team's contribution and formally close out the
    problem-solving process. \\
  \bottomrule
\end{tabular}
\end{minipage}
\vspace{\baselineskip}

For causer attribution, the 8D process is particularly relevant because
its disciplines produce different forms of historical evidence: D2
contains a structured description of the problem, D3 records the interim
containment measures, D4 documents the verified root cause, and D5
through D7 describe the selected, implemented, and preventive corrective
actions.

Completed 8D reports can therefore provide evidence about how comparable
historical incidents were investigated and resolved. However, a
historical 8D report should not automatically determine the causer of a
new defect; its relevance depends on whether the current and historical
cases are comparable in terms of defect, material, supplier, plant, and
process context. For example, an earlier 8D report may document that
comparable door-trim damage was related to supplier packaging. This
provides relevant historical evidence if the current incident involves
the same material, supplier, and incoming inspection point; its relevance
is weak if the defect was first identified only after an internal
assembly activity. Within this thesis, 8D information is consequently
treated as an important historical evidence source rather than as an
automatic attribution rule.

\subsection{Summary}\label{sec:fund_bmw_summary}

This section has described BMW's current quality process and the data it
produces. A defect is first recorded in IPS-Q through MQC and described
using the VFC, whose defect type, defect place, and defect location form
the defect signature introduced in Section~\ref{sec:fund_bmw_vfc}; when
formal investigation is required, this is escalated to a QI, which may
bundle several defect occurrences or, conversely, split identical-looking
defects across several QIs depending on material, supplier, plant, and
process context (Section~\ref{sec:fund_bmw_qi}). The central decision at
this stage is process-level causer attribution, which must be kept
distinct from technical root-cause identification and which the QMT
front-desk process supports through a fixed set of documented business
questions (Section~\ref{sec:fund_bmw_causer}). Today, this decision relies
on a structured but exact-match historical lookup across supplier, part,
and VFC attributes, which is transparent but fails whenever no identical
historical case exists (Section~\ref{sec:fund_bmw_lookup}). Inspection,
sorting, and rework records provide complementary evidence about where a
defect was detected and how it was corrected, but none of these, nor the
unit that performs them, can be equated with causation
(Section~\ref{sec:fund_bmw_isr}). Several roles, from production and
assembly to quality engineers, internal units, and suppliers, contribute
evidence at different stages of this process
(Section~\ref{sec:fund_bmw_roles}), and supplier-related incidents
additionally generate 8D reports, whose eight disciplines yield
further, differently reliable forms of historical evidence depending on
how far an investigation has progressed (Section~\ref{sec:fund_bmw_8d}).

Across all eight subsections, the same limitation recurs: relevant
evidence is distributed across IPS-Q, QCat, Qcube, SRI orders, and 8D
reports, and today it is combined only through exact-match retrieval on a
small set of attributes, rather than through systematic reasoning over
all available historical evidence. This motivates the knowledge-graph-based
approach developed in the remainder of this thesis, which integrates
these heterogeneous data sources into a single structure
(Chapter~\ref{ch:problem}) and supports evidence retrieval, aggregation,
and explanation beyond exact-match lookup (Chapter~\ref{ch:method}).
